\documentclass[aspectratio=169]{beamer}

% Shared Beamer settings for Elements of AIML lectures (UPES).
% Keep this file free of \documentclass to allow reuse via \input.

% Allow lecture sources to use shared theme/assets without copying them into each lecture folder.
% NOTE: These paths are relative to `Unit-xx/Lecture-yy/latex/` (and `Lab/Experiment-xx/latex/`).
\makeatletter
\def\input@path{{../../../_shared/}{../../../_shared/UPES-Beamer-Theme/}}
\makeatother

\usetheme{UPES}
\setbeamertemplate{navigation symbols}{}

\usepackage{amsmath}
\usepackage{amssymb}
\usepackage{booktabs}
\usepackage{graphicx}
\usepackage{hyperref}
\usepackage{xcolor}
\usepackage{listings}

% Code listings (general-purpose).
\lstset{
  basicstyle=\ttfamily\small,
  keywordstyle=\color{blue},
  commentstyle=\color{gray},
  stringstyle=\color{teal},
  showstringspaces=false,
  columns=fullflexible,
  breaklines=true
}

\usepackage{tikz}
\usetikzlibrary{arrows.meta,positioning,shapes.geometric,calc}

% Per-slide source line (references shown on the slide itself).
\newcommand{\slidesources}[1]{%
  \vfill
  {\tiny\color{gray}\raggedright\sloppy\parbox{\linewidth}{Sources: #1}\par}%
}

% Unit-wise source strings for this subject (used by slides/notes).
% Path is relative to the lecture `latex/` folder (the compilation CWD).
% Source strings for Elements of AIML (used by slides + notes).

\newcommand{\AIMLRepoURL}{https://github.com/tali7c/Elements-of-AIML}

\newcommand{\AIMLLabSources}{%
Provost and Fawcett, \textit{Data Science for Business}.;%
McKinney, \textit{Python for Data Analysis}.;%
WEKA Documentation (Waikato Environment for Knowledge Analysis), accessed 2026-02-28.;%
Matplotlib and Seaborn documentation, accessed 2026-02-28.%
}

\newcommand{\AIMLUnitISources}{%
Rich and Knight, \textit{Artificial Intelligence}, McGraw Hill, 2017.;%
Russell and Norvig, \textit{Artificial Intelligence: A Modern Approach} (4e), Ch. 1--2 (Introduction, Intelligent Agents).;%
Poole and Mackworth, \textit{Artificial Intelligence: Foundations of Computational Agents} (2e), selected sections.%
}

\newcommand{\AIMLUnitIISources}{%
Russell and Norvig, \textit{Artificial Intelligence: A Modern Approach} (4e), Ch. 7--9 (Logical Agents, First-Order Logic, Inference).;%
Huth and Ryan, \textit{Logic in Computer Science} (2e), selected sections on propositional and first-order logic.;%
Genesereth and Nilsson, \textit{Logical Foundations of Artificial Intelligence}, selected chapters.%
}

\newcommand{\AIMLUnitIIISources}{%
Mueller and Massaron, \textit{Machine Learning for Dummies}, For Dummies, 2016.;%
Mitchell, \textit{Machine Learning}, selected chapters on learning basics and evaluation.;%
Scikit-learn documentation, accessed 2026-02-28.%
}

\newcommand{\AIMLUnitIVSources}{%
Mueller and Massaron, \textit{Machine Learning for Dummies}, For Dummies, 2016.;%
James, Witten, Hastie, and Tibshirani, \textit{An Introduction to Statistical Learning} (2e), selected chapters.;%
Scikit-learn documentation, accessed 2026-02-28.%
}

\newcommand{\AIMLUnitVSources}{%
NIST, \textit{AI Risk Management Framework (AI RMF 1.0)}, 2023.;%
Barocas, Hardt, and Narayanan, \textit{Fairness and Machine Learning} (online book), selected sections.;%
Knaflic, \textit{Storytelling with Data}, selected chapters on communicating results.%
}



\title[Elements of AIML]{Elements of AIML}
\subtitle{Unit 04 -- Lecture 03: Clustering, Reinforcement Learning \& Semi-Supervised ML}
\author{Tofik Ali}
\institute{School of Computer Science, UPES Dehradun}
\date{\today}

\graphicspath{{../images/}{../../../_shared/}{../../../_shared/UPES-Beamer-Theme/}}

\begin{document}

\UPESCoverFrame
\setcounter{framenumber}{0}

\begin{frame}
  \titlepage
  \vspace{0.3em}
  \begin{center}
    \footnotesize Topic focus: \textbf{unsupervised} and \textbf{interaction-based} learning paradigms\\[0.25em]
    \footnotesize Repository: \texttt{\AIMLRepoURL}
  \end{center}
  \slidesources{\AIMLUnitIVSources}
\end{frame}

\begin{frame}{Quick Links}
  \centering
  \hyperlink{sec:cluster}{\beamerbutton{Clustering}}\hspace{0.6em}
  \hyperlink{sec:kmeans}{\beamerbutton{k-means}}\hspace{0.6em}
  \hyperlink{sec:rl}{\beamerbutton{RL Basics}}\hspace{0.6em}
  \hyperlink{sec:semi}{\beamerbutton{Semi-Supervised}}\hspace{0.6em}
  \hyperlink{sec:summary}{\beamerbutton{Summary}}
\end{frame}

\begin{frame}{Agenda}
  \tableofcontents
\end{frame}

\section{Clustering}
\label{sec:cluster}

\begin{frame}{Learning Outcomes}
  By the end of this lecture, you should be able to:
  \begin{itemize}[<+->]
    \item Explain clustering and where it is used.
    \item Describe the k-means algorithm and its objective.
    \item Explain RL terminology (state, action, reward, policy) at a high level.
    \item Describe semi-supervised learning and when it is helpful.
  \end{itemize}
\end{frame}

\begin{frame}{Abbreviations \& Symbols}
  \begin{itemize}[<+->]
    \item Clustering: $k$ clusters,\; $\mu_j$ centroid of cluster $j$
    \item Objective (inertia): $\sum_i \|x_i-\mu_{c(i)}\|^2$
    \item RL: $s$ state,\; $a$ action,\; $r$ reward,\; $\pi$ policy,\; $\gamma$ discount
    \item $Q(s,a)$: action-value (Q-function)
    \item Semi-supervised: labeled set $L$, unlabeled set $U$, pseudo-labels
  \end{itemize}
\end{frame}

\begin{frame}{What is Clustering?}
  \begin{block}{Clustering}
    Group data points into clusters such that points in the same cluster are more similar.
  \end{block}
  \vspace{0.4em}
  Examples:
  \begin{itemize}[<+->]
    \item customer segmentation
    \item document grouping
    \item anomaly detection (as ``far from clusters'')
  \end{itemize}
\end{frame}

\section{k-means}
\label{sec:kmeans}

\begin{frame}{k-means (Algorithm Sketch)}
  \begin{enumerate}[<+->]
    \item Choose $k$ cluster centers (initialize).
    \item Assign each point to the nearest center.
    \item Update each center as the mean of its assigned points.
    \item Repeat until convergence.
  \end{enumerate}
\end{frame}

\begin{frame}{k-means Objective}
  k-means minimizes within-cluster squared distance:
  \[
    \sum_{i=1}^{n}\left\|x_i - \mu_{c(i)}\right\|^2
  \]
  where $\mu_{c(i)}$ is the center of the cluster assigned to $x_i$.
\end{frame}

\begin{frame}{Mini Worked Example (k-means intuition)}
  Points in 1D: $[1,2,8,9]$, choose $k=2$.
  \begin{itemize}[<+->]
    \item Initialize centers: $\mu_1=2$, $\mu_2=9$
    \item Assign: $\{1,2\}\rightarrow \mu_1$ and $\{8,9\}\rightarrow \mu_2$
    \item Update: $\mu_1=\text{mean}(1,2)=1.5$, $\mu_2=\text{mean}(8,9)=8.5$
  \end{itemize}
  Repeat until assignments stop changing.
\end{frame}

\begin{frame}{Choosing k (Practical)}
  \begin{itemize}[<+->]
    \item Elbow method (plot objective vs $k$)
    \item Domain knowledge
    \item Stability across runs (k-means depends on initialization)
  \end{itemize}
\end{frame}

\begin{frame}{Clustering Evaluation (Intuition)}
  \begin{itemize}[<+->]
    \item Internal metrics (no labels): inertia (k-means objective), silhouette score
    \item Visual check (2D/3D after PCA) can reveal separation
    \item Beware: ``best'' clustering depends on the business goal
  \end{itemize}
\end{frame}

\section{Reinforcement Learning}
\label{sec:rl}

\begin{frame}{Reinforcement Learning (High Level)}
  \begin{itemize}[<+->]
    \item Agent interacts with environment: \textbf{state} $\rightarrow$ \textbf{action} $\rightarrow$ \textbf{reward}
    \item Goal: learn a \textbf{policy} that maximizes cumulative reward
    \item Feedback is delayed and noisy (harder than supervised learning)
  \end{itemize}
\end{frame}

\begin{frame}{Return and Discounting (Key RL Idea)}
  \begin{itemize}[<+->]
    \item Goal is not just immediate reward, but \textbf{total return}:
    \[
      G_t = r_{t+1} + \gamma r_{t+2} + \gamma^2 r_{t+3} + \dots
    \]
    \item $\gamma \in [0,1)$ discounts future rewards.
    \item Higher $\gamma$ values emphasize long-term planning.
  \end{itemize}
\end{frame}

\begin{frame}{Q-Learning (One-line Update)}
  \begin{itemize}[<+->]
    \item Learn an action-value function $Q(s,a)$.
    \item Update rule:
    \[
      Q(s,a) \leftarrow Q(s,a) + \alpha\bigl(r + \gamma \max_{a'}Q(s',a') - Q(s,a)\bigr)
    \]
    \item $\alpha$ is learning rate; $s'$ is next state after taking action $a$ in state $s$.
  \end{itemize}
\end{frame}

\begin{frame}{Toy Example: Grid World}
  \begin{itemize}[<+->]
    \item State: agent position on a grid
    \item Actions: up/down/left/right
    \item Reward: +1 at goal, -1 at trap, small penalty per step
    \item Learn a policy to reach goal quickly
  \end{itemize}
\end{frame}

\section{Semi-supervised learning}
\label{sec:semi}

\begin{frame}{Semi-Supervised Learning (Idea)}
  \begin{itemize}[<+->]
    \item Use a small labeled dataset + large unlabeled dataset.
    \item Useful when labeling is expensive but unlabeled data is abundant.
    \item Example approach: \textbf{self-training} (pseudo-label confident predictions).
  \end{itemize}
  \vspace{0.4em}
  Risks: error reinforcement if pseudo-labels are wrong.
\end{frame}

\begin{frame}{Semi-Supervised Methods (Examples)}
  \begin{itemize}[<+->]
    \item \textbf{Self-training}: add confident pseudo-labels iteratively
    \item \textbf{Co-training}: two ``views'' of features teach each other (when applicable)
    \item \textbf{Graph-based label propagation}: spread labels to neighbors
  \end{itemize}
  \vspace{0.3em}
  \small Use validation to avoid reinforcing mistakes.
\end{frame}

\section{Summary}
\label{sec:summary}

\begin{frame}{Summary}
  \begin{itemize}
    \item Clustering finds structure without labels; k-means is a common baseline.
    \item RL learns from reward signals in interaction.
    \item Semi-supervised learning combines labeled and unlabeled data.
  \end{itemize}
  \slidesources{\AIMLUnitIVSources}
\end{frame}

\begin{frame}{Exit Questions}
  \begin{enumerate}
    \item What is the k-means objective function (in words)?
    \item Why can k-means give different clusters on different runs?
    \item Define state, action, reward, and policy in RL.
    \item Give one benefit and one risk of semi-supervised learning.
  \end{enumerate}
\end{frame}

\end{document}
